\chapter{Review scope and source revisions}
\label{app:review}

This review checks the booklet against public repository source snapshots
retrieved on 21 October 2026, all on \texttt{main}; the pull requests the
earlier review assumed (\repo{Geometry} \#101, \#102, \#116,
\repo{sea-cucumber} \#42) are merged. The whole of Part~II was run end to
end on lxplus at these revisions (\file{tools/dry\_run.sh}), with the
field maps of \repo{field\_maps}. The revisions are \emph{not} a jointly
released software stack. The whole chain is on data
model 0.6.

\section{Checked primary sources}

The revision links open the exact source trees used for the review.
\begin{center}\small
\begin{tabularx}{\linewidth}{l l L}
\toprule Repository & Revision & Main checks \\ \midrule
Geometry & \href{https://github.com/ShipSoft/Geometry/tree/45fbeb6f9dcbdfa0b6116112200cc089f7b784c8}{\texttt{45fbeb6f9d}} & release 0.3.0: coordinates, volume names, envelopes, known overlaps, tasks, \cmd{build\_geometry} options, class names \\
aegir & \href{https://github.com/ShipSoft/aegir/tree/a65eed53d03ce14b68df49d3f922c04c283367d3}{\texttt{a65eed53d0}} & fixed-target source, scoring planes, plane replay, sensitive-volume entries, output fields \\
data-model & \href{https://github.com/ShipSoft/data-model/tree/ba0dda8bb4b43b9d2a2ded3ede57bdb519ae8d12}{\texttt{ba0dda8bb4}} & 0.6 member names, detector ids, backward-read exceptions \\
field\_service & \href{https://github.com/ShipSoft/field_service/tree/6587979624a98713e5c7045b69b4be39f6dc1797}{\texttt{6587979624}} & map format, tools and their options \\
field\_maps & \href{https://github.com/ShipSoft/field_maps/tree/f6234a00d2d5b2b257df410cd7c1459ad3bc236e}{\texttt{f6234a00d2}} & map files, map names, placements \\
geometry\_service & \href{https://github.com/ShipSoft/geometry_service/tree/8e6792570fad3dce7349aeb992e734f860b1356f}{\texttt{8e6792570f}} & 0.6 interface, namespace, CMake target \\
Shannon & \href{https://github.com/ShipSoft/Shannon/tree/7d9f89a47a09c6865124e532c49ef9206661da57}{\texttt{7d9f89a47a}} & input layout, \texttt{pot}, smearing and per-object output \\
sea-cucumber & \href{https://github.com/ShipSoft/sea-cucumber/tree/15ce818e39d8d21d508bd9dc327a706b738f7125}{\texttt{15ce818e39}} & main (with PR \#42): tasks, options, view rules, data model 0.6 \\
Trout & \href{https://github.com/ShipSoft/Trout/tree/d397c33d38108e7beaf4e32d9c46b7151200e886}{\texttt{d397c33d38}} & hit-stream input and reconstruction scope \\
.github & \href{https://github.com/ShipSoft/.github/tree/141170c66caa785af6b5cf30551359dc47e994c1}{\texttt{141170c66c}} & branch rulesets, merge methods and required checks (\cref{sec:pullrequests}) \\
\bottomrule\end{tabularx}\end{center}

Particularly relevant files are \file{aegir/workflows/lib.libsonnet} and
\file{fixed\_target\_production.jsonnet} and \file{plane\_replay.jsonnet}
(which the booklet's two simulation workflows reproduce),
\file{aegir/src/sensitive\_volumes.hpp} (the detector of each pattern),
\file{aegir/docs/geant4\_integration.md} (scoring planes and SD modes),
\file{Shannon/src/digitise\_hits.cpp} (time window and random streams) and
\file{data-model/tests/data/README.md} (schema migration exceptions). Full
commit IDs are in \file{reviewed\_sources.json}; the accompanying
\file{REVIEW.md} records changes and validation limits.

\section{Before the tutorial session}

\begin{enumerate}
\item Freeze compatible package revisions and lock files. Verify all packages
on lxplus with \cmd{pixi cvmfs}: \file{tools/dry\_run.sh} in the organiser
archive runs Part~II end to end and collects timings, disk use and the
numbers the \textbf{TBD} items need.
\item Provide the field maps, their names and coordinate placements. Check
the maps against the exact geometry database used by all consumers.
\item Time stage~1 and choose the number of protons per participant; produce
a larger sample with the same workflow for the analysis.
\item Run the full chain on lxplus, verify schema readback in each
environment, and replace the remaining \textbf{TBD} values with measured
results (including runtime and disk requirements).
\end{enumerate}

\section{Accessibility and validation}

The booklet uses dark blue and neutral tones on white, with dark text and
light box backgrounds. Box titles state their meaning; arrows distinguish
data, services and the future branch by solid, dashed and dotted strokes.
No task or warning requires distinguishing hues. The supplied black SHiP
logo is included without recolouring. Analysis plots use a high-contrast
single-series line; Gaussian fits additionally use a black dashed line.
This supports greyscale use and avoids dependence on red--green perception;
it is not a claim of formal PDF accessibility certification.

Both editions are compiled and visually checked. The analysis scripts are
run in \file{tests/test\_analysis\_helpers.py} against synthetic events with
the data-model 0.6 member names and a stand-in for ROOT. The four workflows
are evaluated with \cmd{jsonnet} in \file{tests/test\_workflows.py}: the
production workflow expands to the same job as \pkg{aegir}'s
\file{fixed\_target\_production.jsonnet}, the replay to \pkg{aegir}'s
\file{plane\_replay.jsonnet} plus the detectors, the field and
energy-deposit hits, the muon gun has the same detectors and field as the
replay, and every volume pattern is checked against the
\repo{Geometry} sources. No Geant4, digitisation or display run is claimed:
the review environment has no SHiP packages, field maps or CVMFS.

For Part~I, the script \file{tests/test\_git\_workflow.sh} runs every git
command sequence of \cref{ch:git,ch:branches,ch:remotes,ch:team} against local bare repositories standing
in for GitHub (forks, an out-of-date pull request, a rebase with and without
conflict, a rebase-only merge queue). The team exercise (skeleton, per-task
checker, sample generator and reference solutions) runs end to end in
\file{tests/test\_teamwork.py} against a numerical stand-in for PyROOT
(\file{tests/fakeroot/}); it has not yet been run with real ROOT. The pixi manifests and the task,
feature and environment commands of \cref{ch:pixi} were checked with
pixi~0.72.0 (the version on CVMFS) and 0.82 on workspaces without packages (the package channels were not reachable from
the review environment, so no environment was solved and no lock file for the
team repository exists yet). \cmd{pixi cvmfs} exists only on CVMFS and has
not been run: the booklet assumes that \cmd{pixi cvmfs <task or command>}
works as \cmd{pixi run} does in any workspace, and that bare \cmd{pixi cvmfs}
opens a shell, as in FairShip. Both workflows of \cref{ch:ci} pass actionlint
1.7.12, and all tutorial Python passes the ruff configuration CI uses. None of
the workflows has run on GitHub.
